\chapter{Alat GRC dan Peta Jalan Tata Kelola Terpadu}

\section*{Tujuan Pembelajaran}
\begin{itemize}
  \item Membandingkan pilihan alat GRC, dari spreadsheet sampai platform terintegrasi, berdasarkan kebutuhan dan kematangan organisasi.
  \item Menyusun peta jalan GRC bertahap yang selaras dengan peta jalan transformasi digital.
  \item Mempertahankan rekomendasi tata kelola, risiko, dan kepatuhan di depan pemangku kepentingan.
\end{itemize}

\section{Pendahuluan}

% TODO
% Cakupan: register risiko, kepatuhan, dan audit tersebar di puluhan spreadsheet yang isinya saling bertentangan, sehingga direksi menerima angka yang berbeda dari tiap unit.

\section{GRC Terpadu}

% TODO
% Cakupan: gagasan satu kontrol untuk banyak kewajiban; pustaka kontrol bersama dan pemetaan antarkerangka; manfaat dan batas pendekatan terpadu.
% Kasus Medikara: kontrol peninjauan hak akses yang sekaligus memenuhi UU PDP, ISO/IEC 27001, dan NIST CSF.

\section{Fungsi dan Pilihan Alat GRC}

% TODO
% Cakupan: fungsi alat GRC (register, alur kerja, pengelolaan bukti, pelaporan); spreadsheet terstruktur, alat sumber terbuka seperti CISO Assistant dan eramba, serta platform komersial; kriteria pemilihan dan biaya total kepemilikan.
% Pendalaman teknis: memasang alat GRC sumber terbuka, memuat kerangka ISO/IEC 27001 dan NIST CSF, mengimpor register semester ini, dan mengenal OSCAL untuk pertukaran data kontrol.

\section{Menyusun Peta Jalan GRC}

% TODO
% Cakupan: gelombang inisiatif, ketergantungan, quick win, anggaran, dan indikator keberhasilan; menempelkan peta jalan GRC pada peta jalan transformasi dan portofolio inisiatif (dibahas rinci di IT40513); urutan yang realistis bagi organisasi menengah.

\section{Manajemen Perubahan dan Budaya}

% TODO
% Cakupan: GRC sebagai perubahan perilaku, bukan hanya dokumen; peran teladan pimpinan; pelatihan dan kesadaran; bentuk resistensi yang umum dan cara menanganinya.

\section{Mempresentasikan Rekomendasi kepada Direksi}

% TODO
% Cakupan: struktur presentasi C4 (masalah, rekomendasi, biaya, risiko bila tidak dijalankan, keputusan yang diminta); menjawab pertanyaan sulit; menyatukan pesan dengan tiga mata kuliah lain.
% Kasus Medikara: presentasi kepada direksi fiktif Medikara tentang peta jalan "Satu Pasien, Satu Rekam Medis".

\section{Aktivitas: Peta Jalan GRC dan Gladi Presentasi C4}

% TODO
% Tujuan aktivitas: memadukan seluruh luaran semester menjadi rencana bertahap yang siap diputuskan direksi.
% Peran dalam tim: semua peran terlibat; pemilik konteks memastikan rekomendasi realistis; analis risiko menjaga keselarasan dengan risk appetite; analis regulasi memeriksa kewajiban yang belum tertangani; analis teknologi menilai kelayakan alat GRC.
% Langkah: kumpulkan luaran Bab 1 sampai 13; susun peta jalan GRC tiga tahun (gelombang, ketergantungan, pemilik, biaya, indikator); pilih pendekatan alat GRC beserta alasan dan trade-off; selaraskan dengan peta jalan IT40513; gladi presentasi C4 dengan format rapat direksi.
% Luaran dan tonggak capstone: peta jalan GRC terpadu, pilihan alat GRC, dan materi presentasi akhir; masuk C4 yang dipresentasikan menjelang UAS.
% Berkas sources yang direncanakan: sources/chapter14/peta-jalan-grc.xlsx, sources/chapter14/kriteria-alat-grc.xlsx, sources/chapter14/templat-presentasi-c4.md

\section{Ringkasan}

% TODO

\section{Latihan}

\begin{enumerate}
  \item Organisasi dengan 300 karyawan ingin membeli platform GRC yang mahal. Argumen apa yang mendukung dan menolak keputusan itu?
  \item Jelaskan gagasan ``satu kontrol untuk banyak kewajiban'' dengan contoh kontrol peninjauan hak akses yang sekaligus memenuhi UU PDP, ISO/IEC 27001, dan NIST CSF.
  \item Susun tiga inisiatif \textit{quick win} GRC untuk Medikara Health pada enam bulan pertama, lalu jelaskan alasan urutannya.
\end{enumerate}
